\documentclass[10pt,a4paper]{amsart}
\usepackage[margin=19mm]{geometry}
\usepackage{amsmath,amssymb,mathtools}
\usepackage[scr=rsfs,cal=euler]{mathalfa}
\usepackage{xfrac}
\usepackage[unicode,hidelinks,bookmarks=false]{hyperref}
\newcommand{\ie}{i.e.\ }
\newcommand{\cf}{cf.\ }
\newcommand{\resp}{respectively\ }
\newcommand{\Subsection}{\subsection}
\providecommand{\nobreakdash}{\nobreak\hbox{-}}
\setlength{\emergencystretch}{2em}
\multlinegap0pt
\usepackage{bm}
\usepackage{bbm}
\usepackage{dsfont}
\usepackage{xspace}
\usepackage{cancel}
\usepackage{mathtools}
\usepackage[shortlabels]{enumitem}
\usepackage{amsthm}
\makeatletter
\@ifundefined{thm}{\newtheorem{thm}{Theorem}}{}
\@ifundefined{cor}{\newtheorem{cor}[thm]{Corollary}}{}
\@ifundefined{prop}{\newtheorem{prop}[thm]{Proposition}}{}
\makeatother
\DeclareMathOperator{\ord}{ord}
\newcommand{\As}{\operatorname{As}}
\newcommand{\GL}[1]{\operatorname{GL}_{#1}}	
\newcommand{\Q}{\mathbb{Q}}
\newcommand{\R}{\mathcal{R}}
\newcommand{\T}{\mathbf{T}}
\newcommand{\bs}{\boldsymbol}
\newcommand{\f}{\mathbf{f}}
\newcommand{\g}{\mathbf{g}}
\title[(Algebraic) $ p $-adic Artin formalism of twisted triple pro]{(Algebraic) $ p $-adic Artin formalism of twisted triple product Galois representations over real quadratic fields}
\author{Bhargab Das, Aprameyo Pal}
\date{Source: arXiv:2503.07542v3 (CC BY 4.0)}
\begin{document}
\maketitle

\noindent\emph{Source and licence.} (Algebraic) $ p $-adic Artin formalism of twisted triple product Galois representations over real quadratic fields. Version arXiv:2503.07542v3 (CC BY 4.0), distributed under the Creative Commons Attribution 4.0 International licence, as recorded in the arXiv licence metadata for this exact version and shown on the abstract page https://arxiv.org/abs/2503.07542v3. Full rights notice: licenses/arxiv-2503.07542-CC-BY-4.0.txt.\par\smallskip

\noindent\textit{Editorial scope.} Counted unit: the Thm whose source label is \texttt{None} in the article (source0.tex lines 1222--1224, source label \texttt{None}), delivered together with the article's own complete original proof of it (source0.tex lines 1225--1272, one \texttt{proof} environment of 48 lines). No mathematics is written, repaired or re-derived: statement and proof are the article's own text line for line. Required context retained uncounted: prop:4.16 (lines 1098--1106); cor:4.17 (lines 1109--1119). Every definition, display and label of those lines is retained unchanged. Omitted from this package and not redistributed: 1--1097, 1107--1108, 1120--1221, 1273--2592. Nothing inside a retained excerpt is omitted or altered: every retained body is delivered exactly as the article's lines are written, so no omission receipt of any kind accompanies this package. Citations: the retained excerpts carry no citation command at all, so no bibliography entry of the article is transcribed and no reference is redistributed. Reference adaptation: none was needed and none was made; every reference inside the delivered unit resolves inside it, so no unresolved reference or citation is produced. Printed slips kept literal: no printed word, symbol, hypothesis or number of the counted statement or of its proof was altered. Layout and class: the article's own class is \texttt{amsart} with packages including babel, inputenc, geometry, amssymb, amsthm, amsfonts, amsmath, mathtools, mathrsfs, stmaryrd, graphicx, upgreek, dsfont, multirow; the wrapper is the lane's pinned \texttt{amsart} editorial wrapper at 10pt on A4, which restates the theorem-like environments the retained text needs and adds the article's own macro definitions that the retained text needs (\texttt{\textbackslash As}, \texttt{\textbackslash GL}, \texttt{\textbackslash Q}, \texttt{\textbackslash R}, \texttt{\textbackslash T}, \texttt{\textbackslash bs}, \texttt{\textbackslash f}, \texttt{\textbackslash g}, \texttt{\textbackslash ord}), with extra wrapper packages (bm, bbm, dsfont, xspace, cancel, mathtools, enumitem). The delivered numbering is the unit's own. Assets: no retained span contains a figure environment, a graphics-inclusion command, a PDF-page inclusion, a table or a TikZ picture; no image file is copied into this package, the package declares no asset dependency and no image file is redistributed with it. Licence: version arXiv:2503.07542v3 of this article is distributed under the Creative Commons Attribution 4.0 International licence, as recorded in the arXiv licence metadata for this exact version and shown on the abstract page https://arxiv.org/abs/2503.07542v3. A full rights notice is delivered at licenses/arxiv-2503.07542-CC-BY-4.0.txt. Characters: every retained excerpt is pure ASCII, so the delivered engine, XeLaTeX with its default Latin Modern text font, reports no missing character.
	\begin{prop}\label{prop:4.16}
		---
		\begin{enumerate}[(1)]
			\item If $\ord_{v}(\frak{N}_\g)=1 $,then $ \rho_\g(I_{v}) $ is infinite.
			\item If $ \#(\rho_\g(I_{v}))= \infty $, then there is a short exact sequence of $ \R_{\g}[G_{F_{v}}] $-modules
			$$ 0 \rightarrow \R_{\g}(1) \otimes \mu \rightarrow T_\g \rightarrow \R_{\g} \otimes \mu \rightarrow 0 $$
			where $ \mu \colon G_{F_v} \rightarrow \{\pm 1\} $ is an unramified quadratic character.
		\end{enumerate}
	\end{prop}

	\begin{cor}
		\label{cor:4.17}
		Let $ \#(\rho_\g(I_{v}))= \infty $ and $ \g  $ has trivial central character. Then the following are equivalent.
		\begin{enumerate}[(1)]
			\item $ H^1(I_{v}, T_\g \widehat{\otimes}_{R_\g} T_\g (\lambda)) $ is free.
			\item $ H^1(I_{v}, T_\g(P) \widehat{\otimes}_{S_{P}} T_\g(\lambda)(P)) $ is free for some arithmetic prime $ P $.
			\item The Tamagawa factor for $ T_\g(P) \widehat{\otimes}_{S_{P}} T_\g(\lambda)(P) $ equals 0 for some arithmetic prime $ P $.
			\item The Tamagawa factor for $ T_\g(\lambda)(P) $ equals 0 for some arithmetic prime $ P $.
		\end{enumerate}
		Moreover, whenever these equivalent conditions are satisfied, then $ H^1(I_{\omega}, T_\g) $ is free for any finite index subgroup $ I_\omega < I_v $ with $ p \nmid [I_v \colon I_\omega] $.
	\end{cor}

	\begin{thm}
		Let the assumptions of the last corollary hold for both families of Hilbert cuspforms $ \f $, $ \g $ over $ F $. Assume also that the $p$-part of the Tamagawa factor for a single member of the Hida family $\f$ and a single member of the family $\g$ is trivial. Then $H^1(I_v , V )$ is free for any $ G_{\Q,\Sigma} $-representation $V \in \{\T_1, \T_2, \T_3\}$ and $ v \nmid p\infty $.
	\end{thm}

	\begin{proof}
		As  $ \T_3 = \T_1 \oplus \T_2 $, it is enough to show the theorem for $ V = \T_3 $. \\
		If $ v \notin \Sigma $ then $ I_v $ act trivially on $ \T_3 $ (converse is also true). Hence $H^1(I_v , \T_3 )$ is free. The following cases will arise when $ v \in \Sigma, \: v \nmid p\infty $. 
		\begin{enumerate}[\text{Case} 1.]
			\item If $ 1 < \#\rho_{\f}(I_v),\#\rho_{\g}(I_v) < \infty $ then $ p \nmid \#\rho_{\f}(I_v),\#\rho_{\g}(I_v) $. Therefore we have $ 1< \#\rho_{\T_3}(I_v) < \infty $ and $ p \nmid \#\rho_{\T_3}(I_v) $ , implies $ H^1(I_v , \T_3) $ is free via corollary 2.10.
			\item Let $ 1 < \#\rho_{\f}(I_v) < \infty $ and $ \#\rho_{\g}(I_v) = \infty $. Therefore we have  $ I_v^{(p)} $ acts trivially on $ \As(T_\g)(\lambda) $ (the character $ \mu$ in Proposition 2.11. is always quadratic)
			$$ \T_3^{I_v^{(p)}} = (\T_3^{I_v^{(p)}})^{I_v^{(p)}/I_v^{(p)}} = (T_\f^{I_v^{(p)}} \otimes \As(T_\g)(\lambda))^{I_v^{(p)}/I_v^{(p)}} = 0 $$ 
			as $T_\f^{I_v^{(p)}} = 0$ by corollary 2.15. In particular $ H^1(I_v , \T_3) =0$. 
			\item Let $ \#\rho_{\f}(I_v) = \infty $ and $ 1 < \#\rho_{\g}(I_v) < \infty $. Then define $ I_\omega \coloneqq \rho_{\g}^{-1}(1) $. As $ p \nmid \#\rho_{\g}(I_v) $, then $ p \nmid [I_v \colon I_\omega] $. Hence $ I_\omega $ acts freely on $ \As(\rho_{\g})(\lambda) $ and the module $ H^1(I_\omega, T_\f) $ is free, and hence $ H^1(I_v,  \T_3) =(H^1(I_\omega, T_\f)\widehat{\otimes}\As(T_\g)(\lambda))^{I_v/I_\omega}  $ is also free. 
			\item Let us consider $\#\rho_{\f}(I_v) = \infty$ and $\#\rho_{\g}(I_v) = \infty$. Since $(\T_3)^{I_v^{(p)}} = (T_{\f})^{I_v^{(p)}} \widehat{\otimes}_{O} T_{\g}\widehat{\otimes}_{O} T_{\g}^\theta$ (as $I_v^{(p)}$ acts trivially on $T_{\g} \widehat{\otimes}_{O}T_{\g}^\theta$ by Proposition \ref{prop:4.16}), without loss of generality we can take 
			$I_v^{(p)}$ acts trivially on $T_{\f}$ and $T_{\g}$ by Proposition \ref{prop:4.16}. Hence we get 
			\[
			\rho_{\spadesuit}(t) = 
			\begin{pmatrix}
				1 & a_{\spadesuit}
				\\
				0 & 1
			\end{pmatrix}  \,,\qquad \spadesuit \in \{\f, \g, \g^\theta\}
			\]
			in $\GL{\R_{\spadesuit}}(T_{\spadesuit}) \cong \GL{2}(\R_{\spadesuit})$. 
			Moreover, via the proof of Corollary \ref{cor:4.17} we have $a_{\spadesuit} \in \R_{\spadesuit}^\times$. Hence the double coset 
			$\GL{8}(\R)\ (\bs{\rho}_3(t) - 1)\ \GL{8}(\R)$ is given by the matrix 
			\[
			\begin{pmatrix}
				0 & 1 & 0&&&&&
				\\
				&  0& 0&1&&&&
				\\
				& & 0&0&1&&&
				\\
				& & &0&0&1&&
				\\
				&  & &&0&0&1&
				\\
				&&&&&0&0&1
				\\
				&& &&&&0&0
				\\
				&  & &&&&&0
			\end{pmatrix}\,. 
			\]
			This gives us $H^1(I_v,\T_3)$ is free, as required. 
			
			\item Suppose $\#\rho_{\f}(I_v) = 0$, then $H^1(I_v, \T_3) = T_\f^{\dagger} \widehat{\otimes}_{O} H^1(I_v, T_{\g}^{\dagger}\widehat{\otimes}_{O} T_{\g^{c}}^{\dagger})$ is free by Corollary \ref{cor:4.17}. 
			If $\#\rho_{\g}(I_v) = 0$, then $H^1(I_v, \T_3) = H^1(I_v, T_\f)\widehat{\otimes}_{O} T_{\g}\widehat{\otimes}_{O} T_{\g}^{\theta}$ is free also by Corollary \ref{cor:4.17}. This completes the proof of our proposition. 
		\end{enumerate}
		
	\end{proof}

\end{document}
